% Part: lambda-calculus
% Chapter: syntax
% Section: de-bruijn

\documentclass[../../../include/open-logic-section]{subfiles}

\begin{document}

\olfileid{lam}{syn}{deb}

\olsection{Indeks De Bruijn}

Ekuivalensi-$\alpha$ iku lumrah banget,
amarga suku-suku kang ekuivalen-$\alpha$
``nduweni teges padha.'' Pancen,
kita bisa menehi sintaksis kanggo
suku lambda kang ora mbedakake
suku-suku kang bisa dikonversi-$\alpha$
siji dadi sijine. Cara kang paling
misuwur ngganti variabel nganggo
\emph{indeks De Bruijn}.

Nalika nulis $\lambd[x][M]$,
kita nyebut kanthi cetha yen $x$
iku parameter fungsi, mula $x$
ing $M$ bisa ngrujuk marang
parameter iku. Ing notasi
indeks De Bruijn, parameter ora
duwe jeneng; rujukan marang
parameter ing awak fungsi ditulis
nganggo cacah tingkat abstraksi
ing antarane. Umpamane,
ing $\lambd[x][\lambd[y][y x]]$,
abstraksi njaba ngiket variabel~$x$,
dene abstraksi njero ngiket
variabel~$y$. Subsuku $y x$
ana ing cakupan abstraksi njero:
ora ana abstraksi ing antarane
$y$ lan pangikete~$\lambd[y]$,
nanging ana siji abstraksi
ing antarane $x$ lan
pangikete~$\lambd[x]$.
Mula $0\, 1$ ditulis kanggo $y x$,
lan kabeh suku ditulis
$\lambd[][\lambd[][01]]$.

\begin{defn}
  Suku De Bruijn ditetepake
  kanthi induktif kaya mangkene:
  \begin{enumerate}
  \item $n$, kanthi $n$ sembarang
    wilangan asli.
  \item $PQ$, kanthi $P$ lan $Q$
    loro-lorone suku De Bruijn.
  \item $\lambd[][N]$, kanthi $N$
    suku De Bruijn.
  \end{enumerate}
\end{defn}

Alih-wujud formal saka suku lambda
lumrah menyang suku kanthi
indeks De Bruijn kaya mangkene:
\begin{defn}
  \begin{align*}
    F_\Gamma(x) &= \Gamma(x) \\
    F_\Gamma(PQ) &= F_\Gamma(P)F_\Gamma(Q) \\
    F_\Gamma(\lambd[x][N]) &= \lambd[][F_{x,\Gamma}(N)]
  \end{align*}
  Ing kene $\Gamma$ iku dhaptar
  variabel kanthi indeks wiwit nol;
  $\Gamma(x)$ nuduhake posisi
  muncule variabel $x$ kang
  paling ngarep ing $\Gamma$.
  Umpamane, yen $\Gamma$ yaiku
  $x,y,z$, mula $\Gamma(x)$ yaiku
  $0$ lan $\Gamma(z)$ yaiku $2$.
  
  $x,\Gamma$ nuduhake dhaptar
  kang dipikolehi kanthi nyelehake
  $x$ ing wiwitane $\Gamma$.
  Nerusake $\Gamma$ ing conto mau,
  $w,\Gamma$ yaiku $w,x,y,z$.
\end{defn}

Suku lambda lumrah bisa dipikolehi
malih saka suku De Bruijn kaya
mangkene:

\begin{defn}
  \begin{align*}
    G_\Gamma(n) &= \Gamma[n] \\
    G_\Gamma(PQ) &= G_\Gamma(P) G_\Gamma(Q) \\
    G_\Gamma(\lambd[][N]) &= \lambd[x][G_{x,\Gamma}(N)]
  \end{align*}
  Ing kene $\Gamma$ maneh-lagi
  dhaptar variabel kanthi indeks
  wiwit nol, lan $\Gamma[n]$
  nuduhake variabel ing posisi~$n$.
  Rumus iki mbutuhake posisi
  kasebut ana ing dhaptar.
  Umpamane, yen $\Gamma$ yaiku
  $x,y,z$, mula $\Gamma[1]$
  yaiku $y$.

  Variabel $x$ ing persamaan
  pungkasan dipilih minangka
  variabel sembarang kang ora
  ana ing $\Gamma$.
\end{defn}

Ing kene sawetara asil diwenehake
tanpa bukti:

\begin{prop}
  Yen $M \aconvone M'$
  lan $\Gamma$ sembarang dhaptar
  kang ngemot $\FV{M}$, mula
  $F_\Gamma(M) \eqs F_\Gamma(M')$.
\end{prop}

\end{document}
